% week_template.tex — skeleton for every week's notes.
% Copy into the week folder as MODULE_WeekNN_Notes.tex and fill the marked slots.
% Compile FROM THE WEEK FOLDER with the shared style on the search path:
%   $env:TEXINPUTS = ".;C:\Users\mryx1\repos\YanaNotesGenerator\templates;"
%   pdflatex -interaction=nonstopmode MODULE_WeekNN_Notes.tex   (run twice)
\documentclass[10pt,a4paper]{article}
\usepackage{uninotes}

% Module code, full module name, week number, week title — drives the header and title.
\notesmodule{MODULE}{Full Module Name}{0}{Week Title}

\begin{document}

\notestitle

% Contents page: HARD RULE, must fit on one page, and is always plain black
% and white. Use \notestoc, NOT bare \tableofcontents, every week: it wraps
% the same call with the link colour forced to black (hyperref patches
% \tableofcontents itself at \begin{document}, so a plain \tableofcontents
% here would render with coloured entries again).
% If it does not fit at default sizing, add \tightertoc on the line above
% \notestoc; if it still does not fit, add
% \setcounter{tocdepth}{1} above it too (sections only, no subsections).
\notestoc
\newpage

% ============================ LECTURE CONTENT ============================
% Sections/subsections in lecture delivery order.
%
% Every displayed equation OUTSIDE worked examples goes in a RED eqnbox
% together with its varlist, in the SAME box (first use only):
%
%   \begin{eqnbox}
%   \begin{equation} e_0 = Blv \end{equation}
%   \begin{varlist}
%     \vitem{e_0}{output voltage}{V}
%     \vitem{B}{magnetic flux density}{T}
%     \vitem{l}{conductor length}{m}
%     \vitem{v}{velocity}{m\,s^{-1}}
%   \end{varlist}
%   \end{eqnbox}
%
% Definition vs equation, pick BOTH when both exist in the source:
%   - Source gives a written definition AND a formula for the same concept:
%     write the prose definition in a defbox, and the formula + varlist in
%     an eqnbox directly after it, both present.
%   - Source gives ONLY a formula, no written definition: derive a one
%     sentence definition ("Aspect Ratio is the ratio of...") and put it as
%     the opening line inside the eqnbox, above the equation, rather than
%     inventing a separate defbox.
% (This eqnbox rule does not apply inside \begin{examplebox} worked
% examples; equations there stay plain \begin{align}/\begin{equation}.)
%
% Worked examples that the lecture itself walks through belong INLINE, in
% an examplebox right where the lecture presents them (do not move them to
% the end). Only worked examples that come from a separate worksheet,
% solutions sheet, or revision deck, or that are not tied to one specific
% point in the lecture, go in the trailing "Worked Examples Q\&A" section.
%
% Other boxes (titles ALWAYS in braces, ALWAYS Title Case after the colon,
% e.g. "Key Result: Dependent Suspension Coupling" not "Key Result:
% dependent suspension coupling"):
%   \begin{defbox}{Definition: ...} ... \end{defbox}
%   \begin{keybox}{Key Result: ...} ... \end{keybox}
%   \begin{warningbox}{Pitfall: ...} ... \end{warningbox}
%
% Cropped slide figures: \notesfig{fig01_name.png}{0.6}{Caption text}

\section{Topic One}

% ============================== KEY SUMMARY ==============================
\section{Key Summary}
% Compact recap of the week: key equations, definitions, takeaways.

% ========================= WORKED EXAMPLES Q\&A ==========================
\section{Worked Examples Q\&A}
% ONLY worked examples from worksheets/solutions/revision decks, or ones not
% tied to a specific point in the lecture flow. (Worked examples the lecture
% itself walks through stay inline in the lecture content above, in
% delivery order, not here.) Q&A format mirroring the official solution
% method:
%
%   \begin{examplebox}{Worked Example: Vibrometer Amplitude Ratio}
%     \question{...problem exactly as given...}
%     \wstep{1}{Explain in plain English what is done and why}
%     \begin{align} ... \end{align}
%     \wstep{2}{...}
%     \result{...value and what it means physically...}
%   \end{examplebox}

% ==================== REVISION QUESTIONS AND ANSWERS =====================
\section{Revision Questions and Answers}
% Every question from the revision slides, none skipped:
%   \begin{revisionbox}{Question 1: ...}
%     \question{...}
%     ...full tutor style answer; numerical problems use \wstep/\result...
%   \end{revisionbox}

\end{document}
